\chapter{\IfLanguageName{dutch}{Stand van zaken}{State of the art}}%
\label{ch:stand-van-zaken}

% Tip: Begin elk hoofdstuk met een paragraaf inleiding die beschrijft hoe
% dit hoofdstuk past binnen het geheel van de graduaatsproef. Geef in het
% bijzonder aan wat de link is met het vorige en volgende hoofdstuk.
%
% Pas na deze inleidende paragraaf komt de eerste sectiehoofding.
%
% Dit hoofdstuk bevat je literatuurstudie. De inhoud gaat verder op de
% inleiding, maar zal het onderwerp van de graduaatsproef diepgaand uitspitten.
% De bedoeling is dat de lezer na lezing van dit hoofdstuk voldoende vertrouwd
% is met het onderzoeksdomein om de rest van het verhaal te kunnen volgen
% \autocite{Pollefliet2011}.
%
% Je verwijst bij elke bewering die je doet, vakterm die je introduceert, enz.
% naar je bronnen. In LaTeX kan dat met \texttt{\textbackslash textcite\{\}}
% of \texttt{\textbackslash autocite\{\}}.

Dit hoofdstuk vat de concepten samen die nodig zijn om de gemaakte keuzes te plaatsen: ticketing, web API's, lagen, frontendstate en beveiliging.

\Needspace{4\baselineskip}
\section{Ticketing als ondersteuning van supportprocessen}

Een ticketingsysteem houdt vragen, antwoorden en bijlagen centraal bij. Binnen IT service management sluit dat aan bij het opvolgen van incidenten en service requests, waarbij structuur en traceerbaarheid belangrijk zijn~\autocite{Axelos2019ITIL4Foundation,ISOIEC20000One2018}. Ook Atlassian benadrukt dat tickets aanvragen samenbrengen, zodat opvolging niet over losse kanalen verspreid raakt~\autocite{AtlassianTicketingSystem2026}. Voor Forcebit Ticketing betekent dit dat de eerste beschrijving als eerste bericht in de conversatie wordt bewaard. Zo heeft elk ticket vanaf het begin dezelfde structuur.

\Needspace{4\baselineskip}
\section{Web API en backendstructuur}

ASP.NET Core biedt controllers, dependency injection, middleware en configuratie voor web API's~\autocite{MicrosoftASPNetCore2026}. Die onderdelen passen bij een backend die HTTP-verzoeken verwerkt, authenticatie controleert en services aanroept. Entity Framework Core wordt gebruikt als Object-Relational Mapper: C\#-entiteiten worden via configuraties en migraties gekoppeld aan een relationele databank~\autocite{MicrosoftEFCore2026}. Dat is bruikbaar voor tickets, berichten, bijlagen en gebruikers omdat deze duidelijke relaties hebben.

\Needspace{4\baselineskip}
\section{Single Responsibility en lagen}

Single Responsibility betekent dat een klasse of module een duidelijke hoofdreden heeft om te veranderen. Robert C. Martin koppelt dit principe aan onderhoudbare software: code blijft beter aanpasbaar wanneer verantwoordelijkheden niet door elkaar lopen~\autocite{Martin2017CleanArchitecture}. Daarom verdeelt een gelaagde structuur de logica. Controllers behandelen HTTP en geven gegevens door via Data Transfer Objects (DTO's), objecten die data over een grens transporteren~\autocite{Fowler2003DataTransferObject}. Services voeren use cases uit: afgebakende gebruikersdoelen of systeemgedrag~\autocite{Jacobson1992UseCaseDriven}. Repositories verzorgen databanktoegang en domeinregels bevatten businessregels zoals toegang en statusovergangen.

\Needspace{4\baselineskip}
\section{Frontend state en herbruikbaarheid}

De frontend is gebouwd met Expo React Native. React Native werkt met componenten, terwijl Expo extra tooling aanbiedt voor ontwikkeling en webondersteuning~\autocite{ReactNative2026,ExpoDocs2026}. React Navigation verzorgt schermnavigatie en webroutes~\autocite{ReactNavigation2026}. In Forcebit Ticketing blijven ruwe API-resultaten de bron van waarheid. Zoekresultaten, filters, tellingen, groepen en paginatie worden daarvan afgeleid, zodat schermen niet elk hun eigen versie van dezelfde data bijhouden.

\Needspace{4\baselineskip}
\section{Beveiliging en validatie}

Omdat tickets klantgegevens en bijlagen kunnen bevatten, mag de frontend niet de enige beveiligingslaag zijn. OWASP benadrukt dat toegangscontrole aan serverzijde moet worden afgedwongen~\autocite{OWASPAccessControl2026}. Forcebit Ticketing gebruikt daarom JWT-authenticatie en rolgebaseerde autorisatie. Een JSON Web Token is een compact tokenformaat waarmee claims, zoals gebruikersidentiteit en rol, worden doorgegeven~\autocite{JwtRFC7519}. Frontendvalidatie helpt de gebruiker, maar backendvalidatie blijft noodzakelijk.

\vspace*{\fill}

\begin{figure}[H]
  \centering
  \begin{tikzpicture}[
    concept/.style={draw=title!70, rounded corners=3pt, fill=title!8, align=center, font=\small\bfseries, text width=3.45cm, minimum height=1.0cm, inner sep=4pt},
    result/.style={draw=title, rounded corners=3pt, fill=title!16, align=center, font=\small\bfseries, text width=3.45cm, minimum height=1.0cm, inner sep=4pt},
    arrow/.style={->, thick, draw=title!70, shorten <=2pt, shorten >=2pt}
  ]
    \node[concept] (ticketing) at (0,2.0) {Ticketing\\{\scriptsize centrale context}};
    \node[concept] (backend) at (-5.0,0) {Gelaagde backend\\{\scriptsize API, services, domain, persistence}};
    \node[concept] (frontend) at (5.0,0) {Frontend state\\{\scriptsize zoeken, filters, paginatie}};
    \node[concept] (security) at (0,-2.0) {Beveiliging\\{\scriptsize JWT, rollen, servercontrole}};
    \node[result] (workflow) at (0,0) {Onderhoudbare\\supportworkflow};

    \draw[arrow] (ticketing.south) -- (workflow.north);
    \draw[arrow] (backend.east) -- (workflow.west);
    \draw[arrow] (frontend.west) -- (workflow.east);
    \draw[arrow] (security.north) -- (workflow.south);
  \end{tikzpicture}
  \caption[Samenvatting stand van zaken]{Visuele samenvatting van de stand van zaken. De literatuurstudie toont hoe ticketing, een gelaagde backend, herbruikbare frontendlogica en beveiliging samen bijdragen aan een onderhoudbare supportworkflow.}
  \label{fig:stand-van-zaken-samenvatting}
\end{figure}
\vspace*{\fill}
